% Uebungsblatt Elektrizitaetslehre (German, first-year physics), with a circuitikz circuit.
% Headings: "N. Aufgabe: Titel" (number first) with "(P Punkte)" at the right margin; parts "a)", "b)", ...
\documentclass[11pt,a4paper]{article}
\usepackage[T1]{fontenc}
\usepackage[utf8]{inputenc}
\usepackage{lmodern}
\usepackage[ngerman]{babel}
\usepackage[margin=2.3cm,top=2.0cm,bottom=2.2cm]{geometry}
\usepackage{amsmath}
\usepackage[locale=DE]{siunitx}
\usepackage[siunitx,european]{circuitikz}
\usepackage{caption}
\usepackage{enumitem}

\setlength{\parindent}{0pt}
\setlist[enumerate,1]{label=\alph*), leftmargin=2em, itemsep=0.35em, topsep=0.3em}
\newcommand{\aufgabe}[3]{\bigskip\par\noindent\textbf{#1. Aufgabe: #2}\hfill(#3 Punkte)\par\nopagebreak\smallskip}

\begin{document}

% ---- Kopf -----------------------------------------------------------------------------
{\small Physik II \textperiodcentered{} Elektrizität und Magnetismus \hfill Sommersemester 2027}\\[-0.4em]
\rule{\linewidth}{0.6pt}\\[0.5em]
{\Large\bfseries Übungsblatt 7: Gleichstrom und Kondensatoren}\\[0.4em]
Besprechung in der Übung am 19.\ November 2026 \hfill Abgabe bis 16.\ November 2026

% ---- 1. Aufgabe -----------------------------------------------------------------------
\aufgabe{1}{Gleichstromkreis}{8}
Die Schaltung in Abbildung~1 wird von einer Gleichspannungsquelle mit $U_0 = \qty{12}{V}$ gespeist.
Es gilt $R_1 = \qty{30}{\ohm}$, $R_2 = \qty{60}{\ohm}$, $R_3 = \qty{120}{\ohm}$ und
$R_4 = \qty{30}{\ohm}$.

\begin{center}
\begin{circuitikz}[scale=1, transform shape]
  \draw (0,0) to[battery1, invert, l_=$U_0$] (0,3)
        to[R, l=$R_1$, i=$I$] (3,3)
        to[short, *-] (6,3)
        to[R, l=$R_3$, i=$I_3$] (6,0)
        to[short, -*] (3,0)
        to[R, l=$R_4$] (0,0);
  \draw (3,3) to[R, l_=$R_2$, i_=$I_2$] (3,0);
  \node[above] at (3,3.1) {A};
  \node[below] at (3,-0.1) {B};
\end{circuitikz}
\captionof{figure}{Gleichstromkreis mit vier Widerständen.}
\end{center}

\begin{enumerate}
  \item Berechnen Sie den Ersatzwiderstand $R_{23}$ der Parallelschaltung von $R_2$ und $R_3$ sowie den
        Gesamtwiderstand $R_\mathrm{ges}$ der Schaltung.
  \item Berechnen Sie den Gesamtstrom $I$ und die Spannung $U_\mathrm{AB}$ zwischen den Knoten A und B.
  \item Stellen Sie die Knotenregel für den Knoten A und die Maschenregel für die Masche aus $U_0$,
        $R_1$, $R_2$ und $R_4$ auf. Berechnen Sie damit die Ströme $I_2$ und $I_3$.
\end{enumerate}

% ---- 2. Aufgabe -----------------------------------------------------------------------
\aufgabe{2}{Leistung}{5}
Betrachten Sie die Schaltung aus der 1.\ Aufgabe.
\begin{enumerate}
  \item Berechnen Sie die in jedem der vier Widerstände umgesetzte elektrische Leistung.
  \item Wie groß ist die von der Quelle abgegebene Leistung? Vergleichen Sie sie mit der Summe der
        Leistungen aus a).
  \item Welche Energie gibt die Quelle in zwei Stunden ab? Geben Sie das Ergebnis in Wh und in kJ an.
\end{enumerate}

% ---- 3. Aufgabe -----------------------------------------------------------------------
\aufgabe{3}{Kondensator}{7}
\begin{enumerate}
  \item Ein Kondensator mit $C = \qty{10}{\micro F}$ wird an die Spannung $U_0 = \qty{12}{V}$ angeschlossen.
        Berechnen Sie die Ladung $Q$ und die gespeicherte Energie $E$.
  \item Die Kondensatoren $C_1 = \qty{10}{\micro F}$ und $C_2 = \qty{40}{\micro F}$ werden in Reihe
        geschaltet und an \qty{50}{V} gelegt. Bestimmen Sie die Gesamtkapazität, die Ladung und die
        Spannung an jedem der beiden Kondensatoren.
  \item Ein Plattenkondensator wird aufgeladen und dann von der Quelle getrennt. Anschließend wird der
        Plattenabstand verdoppelt. Wie ändern sich Kapazität, Spannung und gespeicherte Energie?
        Woher stammt die zusätzliche Energie?
\end{enumerate}

\end{document}
